\documentclass[12pt]{article}

\usepackage[T1]{fontenc}
\usepackage{lmodern}
\usepackage{microtype}
\usepackage[left=1in,right=1in,top=0.88in,bottom=0.88in]{geometry}
\usepackage{setspace}
\usepackage{amsmath,amssymb,mathtools}
\usepackage{mathrsfs}
\usepackage{titlesec}
\usepackage{fancyhdr}
\usepackage{needspace}
\usepackage[hidelinks]{hyperref}

\setstretch{1.12}
\setlength{\parindent}{0pt}
\setlength{\parskip}{0.72\baselineskip}
\setlength{\abovedisplayskip}{15pt plus 3pt minus 2pt}
\setlength{\belowdisplayskip}{15pt plus 3pt minus 2pt}
\setlength{\abovedisplayshortskip}{14pt plus 2pt minus 2pt}
\setlength{\belowdisplayshortskip}{14pt plus 2pt minus 2pt}
\setlength{\jot}{10pt}
\raggedbottom
\clubpenalty=10000
\widowpenalty=10000
\displaywidowpenalty=10000

\titleformat{\section}
  {\large\bfseries}
  {}{0pt}{}
\titlespacing*{\section}{0pt}{1.6\baselineskip}{0.62\baselineskip}

\pagestyle{fancy}
\fancyhf{}
\renewcommand{\headrulewidth}{0pt}
\fancyfoot[C]{\small\thepage}

\newcommand{\R}{\mathbb R}
\newcommand{\I}{\mathcal I}

\begin{document}

\begin{center}
  {\LARGE\bfseries How the Proof Works}\par
  \vspace{0.4\baselineskip}
  {\normalsize Thomas Birnie}\par
\end{center}

\vspace{0.65\baselineskip}

The three-dimensional Navier--Stokes regularity problem asks whether a smooth
solution can lose its smoothness in finite time. Fix a smooth, divergence-free
datum \(u_0\) and a viscosity \(\nu>0\).

Take a vortex tube as the physical picture. As it stretches and thins, the
fluid around its shrinking core turns faster and faster, and the velocity
changes across a smaller distance. The tube cannot lengthen without becoming
narrower.

To see where that concentration sits in the Sobolev scale, rescale the fluid:
\[
  u_\lambda(x,t)=\lambda u(\lambda x,\lambda^2t),
  \qquad
  \|u_\lambda(t)\|_{\dot H^s}
  =
  \lambda^{s-1/2}\|u(\lambda^2t)\|_{\dot H^s}.
\]
Below \(s=\tfrac12\), concentration makes the Sobolev norm smaller. Above
\(s=\tfrac12\), it makes the norm larger. At \(s=\tfrac12\), the power of
\(\lambda\) disappears. The exponent vanishes between the energy level below it
and the first-derivative level above it. That unchanged middle level is the
critical height. At this height, viscosity and nonlinear transport have equal
scaling weight.

The kinetic energy can remain finite while smaller and smaller regions carry
steeper gradients.

In order for a singularity to form, velocity has to blow up. In order for that
to happen, its acceleration would also have to blow up, which means its jerk
has to blow up, and so on, and so on:
\[
  u,
  \qquad
  \partial_tu,
  \qquad
  \partial_t^2u,
  \qquad
  \ldots
\]
If a singularity were forming, what we would expect to see as we move up those
time derivatives is less spatial regularity, not more.

\clearpage
\section*{What the time derivatives reveal}

Write out the first three derivatives:
\[
\begin{aligned}
  \partial_tu
  ={}&
  \nu\Delta u
  -(u\cdot\nabla)u
  -\nabla p,
  \\
  \partial_t^2u
  ={}&
  \nu\Delta\partial_tu
  -(\partial_tu\cdot\nabla)u
  -(u\cdot\nabla)\partial_tu
  -\nabla\partial_tp,
  \\
  \partial_t^3u
  ={}&
  \nu\Delta\partial_t^2u
  -(\partial_t^2u\cdot\nabla)u
  -2(\partial_tu\cdot\nabla)\partial_tu
  \\
  &
  -(u\cdot\nabla)\partial_t^2u
  -\nabla\partial_t^2p.
\end{aligned}
\]
By the third line, the nonlinear term is visibly branching. The binomial form
writes the complete expansion at once:
\[
\begin{aligned}
  \partial_t^{r+1}u
  &+
  \sum_{a=0}^{r}\binom{r}{a}
  \bigl((\partial_t^au)\cdot\nabla\bigr)\partial_t^{r-a}u
  +
  \nabla\partial_t^rp
  \\
  &=
  \nu\Delta\partial_t^ru.
\end{aligned}
\]
Set
\[
  V_r:=\partial_t^ru,
  \qquad
  Q_r:=\partial_t^rp.
\]

Now ask what happened to spatial regularity while those time derivatives were
climbing. Write \(\Lambda:=(-\Delta)^{1/2}\). At time-derivative order \(r\)
and spatial height \(s\), measure \(V_r\) by
\[
  E_{r,s}(t):=\frac12\|\Lambda^sV_r(t)\|_2^2.
\]
Call the complete transport and pressure contribution \(\mathcal P_{r,s}\):
\[
\begin{aligned}
  \mathcal P_{r,s}:={}&
  -\operatorname{Re}
  \left\langle
    \Lambda^sV_r,
    \Lambda^s\!\left[
      \sum_{a=0}^{r}\binom{r}{a}
      (V_a\cdot\nabla)V_{r-a}
      +\nabla Q_r
    \right]
  \right\rangle_{L^2}.
\end{aligned}
\]
Taking the \(H^s\) energy of the differentiated equation gives
\[
  E_{r,s}'=\mathcal P_{r,s}-2\nu E_{r,s+1}.
\]
This is where viscosity enters the argument. The time derivative of the energy
at height \(s\) contains the energy at height \(s+1\), with coefficient
\(2\nu\). Start with the original velocity at the critical height,
\(H:=E_{0,1/2}\). Then
\[
  H'=\mathcal P_{0,1/2}-2\nu E_{0,3/2}.
\]
Each derivative of \(H\) is built from the temporal responses above.
Differentiating \(H=\tfrac12\|\Lambda^{1/2}u\|_2^2\) distributes the time
derivatives between its two copies of \(u\):
\[
  H^{(n)}(t)
  =
  \frac12
  \sum_{a=0}^{n}\binom{n}{a}
  \operatorname{Re}
  \left\langle
    \Lambda^{1/2}V_a(t),
    \Lambda^{1/2}V_{n-a}(t)
  \right\rangle_{L^2}.
\]
Iterating the viscous energy identity \(n\geq1\) times gives
\[
  \frac{d^nH}{dt^n}
  =
  (-2\nu)^nE_{0,n+1/2}
  +
  \sum_{j=0}^{n-1}
  (-2\nu)^j
  \partial_t^{\,n-1-j}\mathcal P_{0,j+1/2}.
\]
\Needspace{7\baselineskip}
The top term is the Sobolev energy at height \(n+\tfrac12\). As temporal order
rises, Sobolev height rises with it. That is the reversal. The singularity
picture led us to expect spatial regularity to fall away as we followed
velocity, acceleration, and jerk. The equation exposes higher spatial energy
instead.

In three dimensions, the Sobolev embedding theorem gives
\[
  s>k+\frac32
  \quad\Longrightarrow\quad
  H^s(\R^3)\hookrightarrow C^k(\R^3),
  \qquad
  \bigcap_{m<\infty}H^m(\R^3)
  =
  H^\infty(\R^3)
  \hookrightarrow
  C^\infty(\R^3).
\]
This is the payoff waiting at the endpoint. If every finite height reaches
\(T_*\), the endpoint lies in \(H^\infty\) and is smooth.

\section*{How the heights reach the endpoint}

The calculation has exposed the heights. The finite rectangles carry them to
\(T_*\).

The index \(r\) tells us which time derivative \(V_r\) we are following, and
the index \(m\) tells us the Sobolev height at which it is measured. The pairs
\((r,m)\) form the mixed grid. Cutting that grid off at a finite temporal depth
\(N\) and a finite spatial height \(M\) gives one rectangle.

For each \(n\leq N\), the adjacent pair
\[
  V_n\in L_t^2H_x^{M+1},
  \qquad
  V_{n+1}=\partial_tV_n\in L_t^2H_x^{M-1}
  \quad\Longrightarrow\quad
  V_n\in C_tH_x^M
\]
carries \(V_n\) continuously through \(T_*\) at the middle height \(H^M\).
Collecting finitely many adjacent pairs, together with their pressure rows
\(Q_n\), gives the pressure-complete rectangle \(\mathcal R_{N,M}[u_0]\).

The rectangle estimate keeps every fixed finite block controlled through
\(T_*\). Larger blocks contain smaller ones, and their shared coordinates agree
because every block comes from the fluid generated by \(u_0\). Taking all
finite blocks together gives \(\I[u_0]\), whose velocity belongs to
\[
  \bigcap_{r,m<\infty}
  H^r(0,T_*;H^m(\R^3)).
\]
Every finite temporal order and every finite spatial height are now present at
once.

Return to the zeroth row \(V_0=u\), denoted
\(\mathscr S_0:=\pi_0\I[u_0]\). For every finite \(m\), the intersection gives
\[
  u\in H^1(0,T_*;H^m(\R^3))
  \hookrightarrow
  C([0,T_*];H^m(\R^3)),
  \qquad m<\infty.
\]
The original velocity now has one endpoint in all of those spaces:
\[
  u_*:=u(T_*)
  \in
  \bigcap_{m<\infty}H^m(\R^3)
  =
  H^\infty(\R^3)
  \hookrightarrow
  C^\infty(\R^3).
\]
The singularity picture led us to expect less spatial regularity as the time
derivatives rose. The intersection reaches \(T_*\) with every finite spatial
height intact.

Local smooth existence from \(u_*\) continues the solution through
\([T_*,T_*+\delta]\). Uniqueness makes that continuation agree with the
solution on \([0,T_*)\). This crosses a supposedly maximal finite time, so
\[
  T_* = \infty.
\]

\section*{Where the estimates come from}

Now choose any classical estimate. For finite \(r,k\), \(2\leq p\leq\infty\),
and
\[
  m>k+\frac32-\frac3p,
\]
the intersection and Sobolev embedding give
\[
  \partial_t^ru\in C_tH_x^m
  \quad\Longrightarrow\quad
  \nabla^k\partial_t^ru
  \in L_t^\infty L_x^p.
\]
That one rule gives, for example,
\[
\begin{aligned}
  \I[u_0]
  &\longrightarrow C_tH_x^1
  \longrightarrow L_t^\infty L_x^6,
  \\
  \I[u_0]
  &\longrightarrow C_tH_x^2
  \longrightarrow L_t^\infty L_x^\infty,
  \\
  \I[u_0]
  &\longrightarrow C_tH_x^3
  \longrightarrow L_t^\infty W_x^{1,\infty}.
\end{aligned}
\]
The \(H^3\) line gives, on any finite interval \([0,T]\),
\[
  \int_0^T\|\nabla u(t)\|_{L^\infty}\,dt
  \leq
  CT
  \sup_{0\leq t\leq T}\|u(t)\|_{H^3}
  <\infty.
\]
Change \(r\), \(k\), or \(p\), and choose the corresponding finite height. The
intersection already contains it.

The familiar continuation estimates are read from the intersection afterward.
They are derived from \(u_0\) through the intersection, not the other way
around.

The one-line version is
\[
\begin{aligned}
  u_0
  &\longrightarrow
  \{(V_n,Q_n)\}_{n\geq0}
  \longrightarrow
  \{\mathcal R_{N,M}[u_0]\}_{N,M<\infty}
  \longrightarrow
  \I[u_0]
  \\
  &\longrightarrow
  \mathscr S_0
  \longrightarrow
  u_*\in H^\infty
  \longrightarrow
  T_*=\infty.
\end{aligned}
\]

\end{document}
